\BellaFrontTitle{Préface}
\origpage{IX--X}
Enseigner, c'est aussi conserver une part de fraîcheur d'âme à force de fréquenter de jeunes élèves. Au fur et à mesure que l'on s'éloigne de l'école, on croit encore la comprendre alors que l'on n'est parfois même plus capable de se souvenir de ses propres problèmes d'alors. C'est ainsi qu'il m'a été donné récemment le « douloureux » privilège de revivre un peu de la dureté oubliée des classes préparatoires par l'intermédiaire de ma fille qui a traversé ce terrible parcours du combattant. J'ai pu observer à cette occasion et en la voyant travailler à quel point l'on peut être déconnecté de la réalité scolaire si l'on ne veut pas accepter de faire l'effort de s'y replonger avec modestie. J'ai été tout particulièrement frappé, après qu'elle ait été enfin reçue à l'école de ses souhaits, de voir qu'elle avait classé ses livres et notes de cours en trois groupes. Le premier ne présentait, selon elle, aucun intérêt alors qu'il s'y trouvait bien des livres que j'aurais moi-même cru essentiels. Le second comprenait des textes considérés comme modérément utiles, bien que marginaux relativement au concours. Le troisième, était composé des ouvrages, qui, pour elle, \emph{permettaient d'intégrer}. Parmi ces derniers textes, il y avait les notes de cours de Bernard Gostiaux, qui avait été son professeur de Mathématiques Spéciales au Lycée Saint Louis.

J'avais apprécié leurs qualités de clarté et de dynamisme, et l'opinion de ma fille, dictée par un raisonnement sans fard, me convainquit, plus encore que mon propre jugement, qu'il fallait absolument les éditer afin qu'elles jouent au mieux leur rôle précieux auprès des autres élèves et \BellaCorrection{candidats aux Grandes Écoles}{candidats au Grandes Ecoles}.

Je pense que Bernard Gostiaux m'en voudrait de faire trop ici son éloge. Pourtant, il faut aussi et un peu passer par là. L'ouvrage de Géométrie Différentielle dont il est l'auteur en collaboration avec Marcel Berger, aujourd'hui édité et réédité dans cette même collection \BellaCorrection{a un succès}{à un succès} qui ne se dément pas. La précision et la rigueur de ce dernier livre expliquent ce résultat plus qu'un long discours. La série de fascicules qui composent le présent ouvrage a une autre qualité, qui, je le pense, justifiera plus encore un succès d'édition. Il s'agit d'un texte \emph{vivant}, au sens que, en le lisant, on se sent immédiatement transporté au cœur d'une salle de classe, devant un professeur qui vous y expose son cours \emph{tel quel}. Ayant une certaine expérience personnelle de l'enseignement, je puis attester de la difficulté de bien rendre par l'écriture un langage parlé. Trop souvent les livres sont composés dans un style académique froid et convenu, rebutant tout lecteur n'en connaissant pas déjà et à l'avance le contenu. De ce fait, combien de parents ne se sont pas battus avec leurs enfants et sans résultat pour tenter de leur imposer d'apprendre un sujet dans les livres de classe avant que la chose leur ait été enseignée ? Le manque de succès de tels efforts est à comparer avec le plaisir de manger un repas chaud et le dégoût qu'inspirent des restes rassis. Rien ne peut remplacer un cours bien fait, car il est porteur de toute l'énergie et l'enthousiasme que le professeur met à transmettre un sujet qui le passionne.

Je tiens à remercier Bernard Gostiaux d'avoir accepté de publier son cours en l'état. En le lisant et le relisant, je me suis mis à regretter tout de bon de ne pas pouvoir retourner sur les bancs de l'école avec 30 ans de moins !

\begin{flushright}\bfseries Paul Deheuvels\end{flushright}
\clearpage
